\chapter{Lời Nhắn Của Thầy Giáo}
\markboth{Lời Nhắn Của Thầy Giáo}{A Teacher's Final Messages}

\section{Nếu Em Đang Mệt, Hãy Nghỉ Một Chút Nhưng Đừng Tự Bỏ Mình}
\begin{truyen}{Nếu Em Đang Mệt, Hãy Nghỉ Một Chút Nhưng Đừng Tự Bỏ Mình}{If You Are Tired, Rest a While but Do Not Abandon Yourself}
\chuhoa{C}{ó những lúc em cần nghỉ thật. Điều đó không làm em yếu đi ngay lập tức.}
\textit{(There are times when you truly need rest. That does not make you weak.)}

Nhưng thầy muốn nhắn một điều: nghỉ để hồi lại khác với bỏ mặc luôn cuộc học của mình.
\textit{(But I want to say this: resting to recover is different from abandoning your learning entirely.)}

Nếu em đang mệt, hãy thở, ngủ, đi bộ, im một chút, xin bớt một chút, rồi quay lại bằng một bước nhỏ hơn.
\textit{(If you are tired, breathe, sleep, walk, stay quiet for a while, ask for a little relief, then return with a smaller step.)}

Đừng vì một giai đoạn kiệt sức mà kết luận em không còn đi được nữa.
\textit{(Do not conclude from one exhausted season that you can no longer go on.)}
\end{truyen}

\begin{baihoc}
Nghỉ để cứu mình là khôn. Bỏ luôn mình vì quá mệt là điều thầy không muốn em làm.
\textit{(Resting to save yourself is wise. Abandoning yourself because you are exhausted is what I do not want you to do.)}
\end{baihoc}

\begin{ghinhoanh}
\textbf{Short English to remember:}

Rest if needed, but do not abandon yourself.

\end{ghinhoanh}

\begin{tuhocnhanh}
1. Nghỉ để hồi lại khác với buông luôn ở đâu? \textit{How is restorative rest different from giving up?}

2. Khi quá mệt, bước quay lại nên nhỏ tới mức nào? \textit{How small should the return step be when you are exhausted?}

3. Câu nào em cần nghe nhất lúc mình kiệt sức? \textit{What sentence do you most need to hear when exhausted?}
\end{tuhocnhanh}
\ngancach

\section{Nếu Em Đang Thua, Điều Quan Trọng Nhất Là Đừng Nói Chuyện Với Mình Quá Tàn Nhẫn}
\begin{truyen}{Nếu Em Đang Thua, Điều Quan Trọng Nhất Là Đừng Nói Chuyện Với Mình Quá Tàn Nhẫn}{If You Are Losing Right Now, Do Not Speak to Yourself Cruelly}
\chuhoa{K}{hi em đang thua, đầu óc rất hay trở thành kẻ đánh em mạnh nhất.}
\textit{(When you are losing, your mind often becomes the one that hits you hardest.)}

Em có thể tự gọi mình bằng những lời mà nếu người khác nói ra, em sẽ thấy rất ác.
\textit{(You may call yourself names that would sound cruel if another person said them aloud.)}

Thầy không bảo em tự vuốt ve mình giả tạo. Thầy chỉ muốn em sửa mình bằng giọng đủ người.
\textit{(I am not telling you to flatter yourself falsely. I only want you to correct yourself in a voice that is still human.)}

Vì con người bị chửi mãi bởi chính mình sẽ rất khó đứng dậy lâu dài.
\textit{(Because a person constantly cursed by their own mind struggles to rise for long.)}
\end{truyen}

\begin{baihoc}
Tự sửa là cần. Tự hủy mình thì không giúp em đi xa hơn.
\textit{(Self-correction is necessary. Self-destruction will not take you farther.)}
\end{baihoc}

\begin{ghinhoanh}
\textbf{Short English to remember:}

Correct yourself without destroying yourself.

\end{ghinhoanh}

\begin{tuhocnhanh}
1. Vì sao lúc thua mình hay tự nói quá nặng? \textit{Why do we speak harshly to ourselves when losing?}

2. Giọng “đủ người” với bản thân là giọng như thế nào? \textit{What does a humane inner voice sound like?}

3. Em có thể thay một câu tự chửi bằng câu nào? \textit{What sentence could replace self-insult?}
\end{tuhocnhanh}
\ngancach

\section{Đừng Lấy Một Kỳ Thi Làm Toàn Bộ Số Phận Của Em}
\begin{truyen}{Đừng Lấy Một Kỳ Thi Làm Toàn Bộ Số Phận Của Em}{Do Not Turn One Exam into Your Entire Fate}
\chuhoa{K}{ỳ thi quan trọng, thầy không nói dối em điều đó.}
\textit{(Exams matter. I will not lie to you about that.)}

Nhưng một kỳ thi không đủ quyền để nói hết tương lai em, càng không đủ quyền để định nghĩa toàn bộ giá trị em.
\textit{(But one exam does not have enough authority to speak for your whole future, much less define your total worth.)}

Nếu làm tốt, hãy vui nhưng đừng kiêu. Nếu làm chưa tốt, hãy đau vừa đủ rồi nhìn tiếp đường khác, cách khác, nhịp khác.
\textit{(If you do well, be glad but do not become arrogant. If you do poorly, grieve enough, then keep looking for another path, another method, another rhythm.)}

Đời rộng hơn một cột mốc, dù ở thời điểm đó em có thể chưa tin nổi điều này.
\textit{(Life is wider than one milestone, even if at the moment you may find that hard to believe.)}
\end{truyen}

\begin{baihoc}
Hãy nghiêm túc với kỳ thi, nhưng đừng quỳ gối trước nó như trước một vị thần quyết định toàn bộ đời mình.
\textit{(Take exams seriously, but do not kneel before them as though they were gods deciding your entire life.)}
\end{baihoc}

\begin{ghinhoanh}
\textbf{Short English to remember:}

Take exams seriously, but do not worship them.

\end{ghinhoanh}

\begin{tuhocnhanh}
1. Vì sao phải nói thật rằng kỳ thi quan trọng? \textit{Why must we honestly admit that exams matter?}

2. Nhưng vì sao cũng không được tuyệt đối hóa nó? \textit{Why must we also avoid turning it into everything?}

3. Sau một kỳ thi chưa tốt, điều gì nên làm tiếp? \textit{What should come next after a poor exam?}
\end{tuhocnhanh}
\ngancach

\section{Học Chậm Không Có Nghĩa Là Em Đi Không Tới}
\begin{truyen}{Học Chậm Không Có Nghĩa Là Em Đi Không Tới}{Learning Slowly Does Not Mean You Cannot Arrive}
\chuhoa{T}{hầy đã gặp nhiều học sinh nở muộn. Lúc đầu các em đi chậm, nhưng không phải vì thế mà các em không đi tới.}
\textit{(I have met many late bloomers. They began slowly, but that did not mean they never arrived.)}

Có em phải hiểu theo cách khác, phải luyện lâu hơn, phải có người chỉ bình tĩnh hơn.
\textit{(Some needed a different entry point, more practice, or a calmer guide.)}

Đừng vì mình chậm hơn vài người ở đoạn này mà tự kết án cả đường dài.
\textit{(Do not condemn your whole long road just because you are slower than a few people in this stretch.)}

Điều thầy mong là em giữ được việc học đủ lâu để năng lực mình có dịp nở ra.
\textit{(What I hope is that you keep learning long enough for your ability to have a chance to bloom.)}
\end{truyen}

\begin{baihoc}
Nhịp chậm không phải dấu chấm hết. Nó chỉ đòi em bền hơn và tử tế hơn với chính mình.
\textit{(A slower pace is not the end. It simply asks for more endurance and more kindness toward yourself.)}
\end{baihoc}

\begin{ghinhoanh}
\textbf{Short English to remember:}

Slow pace is not the same as no future.

\end{ghinhoanh}

\begin{tuhocnhanh}
1. Vì sao người học chậm dễ tự kết án mình? \textit{Why do slower learners often condemn themselves?}

2. Thầy cô cần giúp người học chậm điều gì nhất? \textit{What do slower learners most need from teachers?}

3. Em cần giữ điều gì để năng lực có thời gian nở ra? \textit{What must you keep so your ability has time to bloom?}
\end{tuhocnhanh}
\ngancach

\section{Nếu Em Chưa Biết Mình Hợp Gì, Cứ Tiếp Tục Lớn Lên Đàng Hoàng Trước Đã}
\begin{truyen}{Nếu Em Chưa Biết Mình Hợp Gì, Cứ Tiếp Tục Lớn Lên Đàng Hoàng Trước Đã}{If You Still Do Not Know What Fits You, Keep Growing Well First}
\chuhoa{K}{hông phải học sinh nào cũng sớm biết mình hợp nghề gì, thích đường nào.}
\textit{(Not every student knows early which career fits or which path they prefer.)}

Nếu em chưa rõ, đừng hoảng quá. Trong lúc chưa rõ, vẫn có những việc rất đáng làm: học cho chắc nền, sống tử tế, tập kỷ luật, tập chịu trách nhiệm, tập giao tiếp cho rõ.
\textit{(If you are still unclear, do not panic. While waiting for clarity, there are still valuable things to do: strengthen your foundations, live decently, practice discipline, responsibility, and clear communication.)}

Những thứ đó không phí, dù sau này em đi đường nào.
\textit{(Those things are never wasted, whatever path you later choose.)}

Nhiều khi đường đi rõ hơn sau khi con người đã lớn lên đủ đàng hoàng để thấy nó.
\textit{(Sometimes the path becomes clearer only after a person has grown into someone steady enough to see it.)}
\end{truyen}

\begin{baihoc}
Chưa rõ hướng không có nghĩa là đứng yên. Em vẫn có thể lớn lên đúng cách trong lúc đang tìm đường.
\textit{(Lacking direction does not mean standing still. You can still grow the right way while searching.)}
\end{baihoc}

\begin{ghinhoanh}
\textbf{Short English to remember:}

You can still grow well while your direction is still unclear.

\end{ghinhoanh}

\begin{tuhocnhanh}
1. Khi chưa rõ hướng, những việc nào vẫn đáng làm? \textit{What is still worth doing when direction is unclear?}

2. Vì sao những nền tảng này không phí? \textit{Why are these foundations never wasted?}

3. Em có thể lớn lên điều gì ngay từ bây giờ? \textit{What can you start growing right now?}
\end{tuhocnhanh}
\ngancach

\section{Người Giỏi Thật Không Nhất Thiết Là Người Làm Em Sợ}
\begin{truyen}{Người Giỏi Thật Không Nhất Thiết Là Người Làm Em Sợ}{A Truly Strong Person Is Not Necessarily One Who Makes Others Afraid}
\chuhoa{N}{hiều học sinh lớn lên với cảm giác rằng người giỏi thì phải sắc lạnh, phải hơn người, phải làm người khác nể vì sợ.}
\textit{(Many students grow up with the feeling that strong people must be sharp, superior, and feared.)}

Thầy không tin như vậy.
\textit{(I do not believe that.)}

Người giỏi thật có thể rất vững mà không cần chèn ép ai. Có thể hiểu nhiều mà không làm người khác thấy bé đi.
\textit{(A truly strong person can be steady without crushing anyone. They can know much without making others feel smaller.)}

Nếu việc học của em làm em khinh người hơn, có lẽ em đang giỏi lên ở một phần mà hỏng đi ở phần khác.
\textit{(If learning makes you despise others more, you may be getting stronger in one part while becoming damaged in another.)}
\end{truyen}

\begin{baihoc}
Sức mạnh đẹp là sức mạnh không cần làm người khác co rúm để chứng minh chính nó.
\textit{(Beautiful strength does not need others to shrink in order to prove itself.)}
\end{baihoc}

\begin{ghinhoanh}
\textbf{Short English to remember:}

Real strength does not need to humiliate others.

\end{ghinhoanh}

\begin{tuhocnhanh}
1. Vì sao nhiều người nhầm giỏi với đáng sợ? \textit{Why do many people confuse strength with intimidation?}

2. Người giỏi thật có thể cư xử khác ra sao? \textit{How can a truly strong person behave differently?}

3. Việc học đang làm em dịu hơn hay khinh người hơn? \textit{Is learning making you gentler or more contemptuous?}
\end{tuhocnhanh}
\ngancach

\section{Nếu Không Có Ai Tin Em Lúc Này, Em Vẫn Phải Giữ Lấy Một Chút Tin Vào Mình}
\begin{truyen}{Nếu Không Có Ai Tin Em Lúc Này, Em Vẫn Phải Giữ Lấy Một Chút Tin Vào Mình}{If No One Believes in You Right Now, Keep at Least a Small Piece of Belief in Yourself}
\chuhoa{C}{ó những giai đoạn học sinh rất cô đơn. Không được khen, không được hiểu, không được ai nghĩ mình làm nên chuyện.}
\textit{(There are seasons when students feel deeply alone. Not praised, not understood, not believed in.)}

Nếu em đang ở đoạn đó, thầy biết rất mệt.
\textit{(If you are in that season, I know it is exhausting.)}

Nhưng thầy mong em giữ lại một chút rất nhỏ thôi: khả năng mình chưa xong, mình chưa phải bản cuối cùng.
\textit{(But I hope you keep one very small thing: the sense that you are not finished yet, not the final version of yourself.)}

Chỉ cần một chút tin còn sống, con người vẫn còn cửa để đi tiếp.
\textit{(As long as a small piece of belief is still alive, there is still a way forward.)}
\end{truyen}

\begin{baihoc}
Không ai tin em lúc này là điều rất đau. Nhưng đừng cộng thêm nỗi đau đó bằng cách tự rút luôn niềm tin cuối cùng khỏi mình.
\textit{(It is painful when no one believes in you. But do not add to that pain by removing the final bit of belief from yourself.)}
\end{baihoc}

\begin{ghinhoanh}
\textbf{Short English to remember:}

You may not be believed in yet, but you are not finished yet.

\end{ghinhoanh}

\begin{tuhocnhanh}
1. Vì sao mất cảm giác được tin rất đau? \textit{Why is losing the sense of being believed in so painful?}

2. “Em chưa phải bản cuối cùng” có thể cứu học sinh ở điểm nào? \textit{How can “you are not your final version yet” help a student?}

3. Một chút niềm tin tối thiểu có thể được giữ bằng cách nào? \textit{How can a minimal self-belief be kept alive?}
\end{tuhocnhanh}
\ngancach

\section{Điều Thầy Mong Nhất Không Phải Là Em Luôn Đứng Đầu}
\begin{truyen}{Điều Thầy Mong Nhất Không Phải Là Em Luôn Đứng Đầu}{What I Hope Most Is Not That You Always Rank First}
\chuhoa{N}{ếu hỏi thầy mong gì nhất ở học trò, câu trả lời không phải là: lúc nào cũng đứng đầu.}
\textit{(If you ask what I want most from my students, the answer is not: to always rank first.)}

Thầy mong em giữ được lòng tử tế, biết học tiếp, biết chịu trách nhiệm, biết sai mà sửa, biết mạnh mà không ác, biết yếu mà không tự bỏ mình.
\textit{(I hope you keep kindness, keep learning, accept responsibility, know how to repair mistakes, be strong without cruelty, and be weak without abandoning yourself.)}

Nếu em có những điều đó, dù đường đời em đi quanh co hơn dự tính, em vẫn có cơ hội trở thành một người đáng quý.
\textit{(If you have those things, then even if life bends more than expected, you still have a chance to become someone deeply worthwhile.)}

Đứng đầu là một vị trí. Làm người cho ra người là một chuyện lớn hơn nhiều.
\textit{(Ranking first is a position. Becoming a good human being is something much larger.)}
\end{truyen}

\begin{baihoc}
Thành tích đẹp là điều đáng mừng. Nhưng điều thầy mong giữ lâu hơn là phẩm chất sống của em.
\textit{(Good achievement is worth celebrating. But what I hope lasts longer is the quality of your character.)}
\end{baihoc}

\begin{ghinhoanh}
\textbf{Short English to remember:}

First place is a position; character is a life.

\end{ghinhoanh}

\begin{tuhocnhanh}
1. Vì sao đứng đầu không thể là điều quan trọng nhất? \textit{Why can first place not be the most important thing?}

2. Những phẩm chất nào thầy mong học trò giữ được? \textit{What qualities does a teacher hope students keep?}

3. Em muốn người khác nhớ mình vì điều gì hơn cả? \textit{What do you most want to be remembered for?}
\end{tuhocnhanh}
\ngancach

\section{Và Nếu Có Một Câu Cuối Cùng Thầy Muốn Nói}
\begin{truyen}{Và Nếu Có Một Câu Cuối Cùng Thầy Muốn Nói}{And If There Is One Final Sentence I Want to Leave You}
\chuhoa{N}{ếu có một câu cuối cùng, thầy muốn nói thế này: em không cần phải hoàn hảo mới xứng đáng được đi tiếp.}
\textit{(If there is one final sentence, it is this: you do not need to be perfect to deserve moving forward.)}

Em chỉ cần còn biết học, còn biết sửa, còn biết giữ lấy phần người trong mình, thì cuộc đời vẫn chưa đóng cửa với em.
\textit{(You only need to keep learning, keep repairing, and keep the human part of yourself alive, and life is not closed to you.)}

Sẽ có lúc em mệt, sai, chậm, lạc hướng, hoặc phải bắt đầu lại.
\textit{(There will be times when you are tired, wrong, slow, lost, or forced to begin again.)}

Nhưng xin đừng vì những lúc đó mà kết luận quá sớm về toàn bộ con người mình.
\textit{(But please do not use those moments to make a final judgment about your whole self too early.)}
\end{truyen}

\begin{baihoc}
Còn học được, còn sửa được, còn giữ được phần người, là còn rất nhiều hy vọng.
\textit{(If you can still learn, still repair, and still keep your humanity, then much hope remains.)}
\end{baihoc}

\begin{ghinhoanh}
\textbf{Short English to remember:}

You do not need perfection to deserve another step.

\end{ghinhoanh}

\begin{tuhocnhanh}
1. Vì sao câu này có thể quan trọng với học sinh? \textit{Why can this sentence matter so much to students?}

2. Ba điều “còn học, còn sửa, còn giữ phần người” gợi cho em điều gì? \textit{What do the three ideas “still learn, still repair, still keep your humanity” mean to you?}

3. Nếu phải giữ lại một câu trong cả quyển này, em giữ câu nào? \textit{If you could keep one sentence from this whole book, which would it be?}
\end{tuhocnhanh}
\ngancach

\section{Thầy Mong Em Đi Tiếp Bằng Một Cách Tử Tế Với Chính Mình Và Với Người Khác}
\begin{truyen}{Thầy Mong Em Đi Tiếp Bằng Một Cách Tử Tế Với Chính Mình Và Với Người Khác}{I Hope You Continue with Decency Toward Yourself and Others}
\chuhoa{C}{uộc đời sau này có thể làm em va nhiều hơn những gì trường học từng chuẩn bị.}
\textit{(Life later may hit you with more than school ever prepared you for.)}

Thầy không mong em đi qua mọi thứ mà không đau.
\textit{(I do not expect you to pass through everything without pain.)}

Thầy chỉ mong khi đi tiếp, em đừng học cách mạnh lên bằng cách độc hơn với người khác hay tàn nhẫn hơn với chính mình.
\textit{(I only hope that as you continue, you do not learn to become strong by becoming more cruel to others or harsher toward yourself.)}

Nếu em giữ được sự tử tế đó, rồi vẫn học, vẫn làm việc, vẫn sửa mình, em đã đi rất đúng đường rồi.
\textit{(If you keep that decency and still continue learning, working, and repairing yourself, then you are already on a very right path.)}
\end{truyen}

\begin{baihoc}
Đi tiếp cho đúng không chỉ là đi tới. Nó còn là đi mà không đánh mất phần tốt của mình.
\textit{(To continue rightly is not only to move forward. It is to move forward without losing the good part of yourself.)}
\end{baihoc}

\begin{ghinhoanh}
\textbf{Short English to remember:}

Go forward without losing the good part of yourself.

\end{ghinhoanh}

\begin{tuhocnhanh}
1. Vì sao sức mạnh không nên được xây bằng sự tàn nhẫn? \textit{Why should strength not be built on cruelty?}

2. “Đi tiếp cho đúng” gợi cho em điều gì? \textit{What does “continue in the right way” mean to you?}

3. Em muốn giữ phần tốt nào của mình lâu nhất? \textit{What good part of yourself do you most want to keep?}
\end{tuhocnhanh}
\ngancach
